\documentclass[12pt]{exam}
\usepackage[utf8]{inputenc}
\usepackage[T1]{fontenc}
\usepackage{wrapfig}
\usepackage{amsmath}
\usepackage{amssymb}
\usepackage{hyperref}
\usepackage{graphicx}
\usepackage[norsk]{babel}
\usepackage{tikz}
\usepackage{tkz-euclide}
\usepackage{float}
%\renewcommand{\headrulewidth}{0pt}
%\fancyhead[L]{}
%\fancyhead[R]{
%\includegraphics[width=4cm]{./img/logo.jpg}}
\begin{document}

\pagestyle{headandfoot}
\runningheadrule
\firstpageheader{}{}{}
\firstpagefooter{}{}{}
\runningheader{Eksamenskurs i matematikk for ungdomskolen 2018}{}{\includegraphics[width=0.2\textwidth]{./img/logo.jpg}}
\runningfooter{}{}{}



{  \centering

	{\scshape\Huge Eksamenskurs i matematikk \par}
	\vspace{1cm}
	{\Large Ungdomskolen \today \par}
	\vspace{1.5cm}
        \begin{center}
          { \large
          \begin{tabular}{c | c}
            Tid & Hva \\ \hline
            16:15-17:30 & Del 1 : Linære funksjoner\\ \hline
            17:30-17:45 & Pause \\\hline
            17:45-19:00 & Del 2 : Linære likninger\\\hline
          \end{tabular}
          }
          \vspace{1.5cm}
          \begin{figure}[h!]
            \centering
            \includegraphics[scale = 0.7]{./img/forside.png}                      
          \end{figure}

        \end{center}

}

\newpage

\qformat{\underline{Spørsmål \thequestion} \hfill}
\section*{Linære likninger}
% Intro
\[
y = a\cdot x + b
\]
Vi kaller da:
\begin{itemize}
\item $a$ stigningstall
\item $b$ kontantledd
\item $x$ variabel
\end{itemize}
Vi oppgir punkter i $(x,y)$ koordinater.

% Oppgave 1
\begin{questions}
\question Sett punktene $(-1,4)-(-2,-5)-(3,1)-(4,9)-(2,-3)$ inn i koordinatsystemet.
\end{questions}
\begin{center}
\begin{tikzpicture}[scale = 0.7,every node/.style={scale=0.6}]
\tkzInit[xmin=-10,xmax=10,ymin=-10,ymax=10]
\tkzGrid[color=gray]
\tkzAxeXY[very thick]
\tkzGrid
\end{tikzpicture}  
\end{center}

% Oppgave 2




\end{document}





